\chapter{Classical special integrals}
\label{module:M03.18}
\SourceRef{Zorich II, section 17.3; Rudin, section 8, Gamma-function part;
Olver, chapter 2 (real positive arguments only)}

\KEntry{K03.18.beta_definition}{The Beta integral and its domain}
\begin{definition}
For $p,q>0$ put
\[
 B(p,q)=\int_0^1 t^{p-1}(1-t)^{q-1}\dd t.
\]
This is an improper integral, with the two endpoints checked independently.
\end{definition}
\begin{proof}[Existence and symmetry]
On $(0,1/2]$ the second factor is bounded above and below by positive constants,
so convergence there follows from $\int_0^{1/2}t^{p-1}\dd t<\infty$.
On $[1/2,1)$ use the same argument with $(1-t)^{q-1}$. Both tests require the
indicated positive parameter; no linked truncation is being used. The integrand
is positive, so $B(p,q)>0$. Substitution $u=1-t$ first on regular compact
subintervals and then at each endpoint gives $B(p,q)=B(q,p)$.
\end{proof}

\KEntry{K03.18.beta_recurrence}{Beta recurrences}
\begin{proposition}
For $p,q>0$,
\[
 B(p+1,q)=\frac p{p+q}B(p,q),\qquad
 B(p,q+1)=\frac q{p+q}B(p,q),\qquad B(p,1)=\frac1p.
\]
Consequently, for integers $n\geq0$,
\begin{equation}\label{eq:beta-integer}
 B(s,n+1)=\frac{n!}{s(s+1)\cdots(s+n)}\quad(s>0).
\end{equation}
\end{proposition}
\begin{proof}
Adding the two integrands gives $B(p+1,q)+B(p,q+1)=B(p,q)$.
The derivative of $t^p(1-t)^q$ is
$pt^{p-1}(1-t)^q-qt^p(1-t)^{q-1}$. Integrating on $[\eta,1-\delta]$ and sending
$\eta,\delta$ independently to zero gives
$pB(p,q+1)=qB(p+1,q)$, since the boundary values tend to zero and both
derivative terms are integrable. Solving these two linear equations gives the
recurrences. Direct evaluation of $\int_0^1t^{p-1}\dd t$ gives $1/p$.
Iterate the recurrence in the second parameter to obtain
\eqref{eq:beta-integer}; this is finite induction, not a Gamma identity.
\end{proof}

\KEntry{K03.18.gamma_definition}{The real Gamma integral}
\begin{definition}
For $s>0$ define
\[
 \Gamma(s)=\int_0^\infty t^{s-1}e^{-t}\dd t.
\]
\end{definition}
\begin{proof}[Existence and positivity]
On $(0,1]$ the integrand is bounded by $t^{s-1}$, whose improper integral is
$1/s$. At infinity every fixed power is eventually dominated by an exponential:
choose an integer $m>s$ and use the exponential series, or
\KRef{K03.02.growth_scale}, to bound $t^{s-1}e^{-t}\leq C e^{-t/2}$ for
$t\geq1$. The latter function has a finite integral. The two positive improper
pieces therefore converge independently. The integrand is strictly positive on
every nonempty compact subinterval of $(0,\infty)$, so $\Gamma(s)>0$.
\end{proof}

\KEntry{K03.18.gamma_recurrence}{Recurrence and factorial values}
\begin{theorem}
For $s>0$, $\Gamma(s+1)=s\Gamma(s)$; moreover $\Gamma(1)=1$ and
$\Gamma(n+1)=n!$ for every integer $n\geq0$.
\end{theorem}
\begin{proof}
On $[\eta,R]$ integrate $t^se^{-t}$ by parts:
\[
 \int_\eta^R t^se^{-t}\dd t
 =[-t^se^{-t}]_\eta^R+s\int_\eta^R t^{s-1}e^{-t}\dd t.
\]
The upper boundary tends to zero by exponential domination; the lower boundary
tends to zero since $s>0$. Both integrals exist by the preceding entry, so the
recurrence follows from separate endpoint limits. The elementary integral of
$e^{-t}$ is $1$ on $(0,\infty)$. Induction now proves the factorial formula,
including $0!=1$.
\end{proof}

\KEntry{K03.18.gamma_euler_limit}{Euler's one-variable limit}
\begin{theorem}
For every $s>0$,
\begin{equation}\label{eq:gamma-euler}
 \Gamma(s)=\lim_{n\to\infty}n^sB(s,n+1)
 =\lim_{n\to\infty}\frac{n!n^s}{s(s+1)\cdots(s+n)}.
\end{equation}
\end{theorem}
\begin{proof}
Substitute $u=nt$ in the Beta integral:
\[
 n^sB(s,n+1)=\int_0^n u^{s-1}(1-u/n)^n\dd u.
\]
For $0\leq u\leq n$ the logarithmic inequality $\log(1-v)\leq-v$ gives
$(1-u/n)^n\leq e^{-u}$; the value at $u=n$ is interpreted by continuity.
On every fixed $[\eta,R]\subset(0,\infty)$,
$n\log(1-u/n)=-u+O(R^2/n)$ uniformly, so the integrands tend uniformly to
$u^{s-1}e^{-u}$. This uniform logarithmic estimate follows, for example, by
integrating $1/(1-v)-1$ on $0\leq v\leq R/n$; its absolute integral is at most
$2(R/n)^2$ for $n\geq2R$.

Here is the full improper-limit argument. On $(0,\eta)$ both nonnegative
integrands are bounded by $u^{s-1}$, so their integrals are at most $\eta^s/s$.
For $u\geq R$ their integrals are bounded by the tail of
$u^{s-1}e^{-u}$, uniformly in $n$, with the first integrand set to zero for
$u>n$. Given $\eps>0$, choose $\eta$ small and $R$ large to make these endpoint
contributions below $\eps$, and then use uniform convergence on $[\eta,R]$.
This proves the limit with ordinary proper integrals and tail estimates.
Equation \eqref{eq:beta-integer} supplies the product formula. No
multidimensional integration theorem is used.
\end{proof}

\KEntry{K03.18.gamma_regular}{Differentiation of Gamma}
\begin{theorem}
The function $\Gamma$ is $C^\infty$ on $(0,\infty)$ and, for $k\geq0$,
\[
 \Gamma^{(k)}(s)=\int_0^\infty t^{s-1}(\log t)^k e^{-t}\dd t.
\]
\end{theorem}
\begin{proof}
Fix $0<\alpha\leq s\leq\beta<\infty$. On $(0,1]$ the absolute value of the
$k$th parameter derivative is at most
$t^{\alpha-1}|\log t|^k$. Substitution $t=e^{-v}$ turns its integral into
$\int_0^\infty e^{-\alpha v}v^k\dd v$, which converges by the Gamma
existence estimate (or repeated integration by parts). On $[1,\infty)$ use
$t^{\beta-1}(\log t)^ke^{-t}\leq C_k e^{-t/2}$. These two integrable bounds
also apply to the next derivative. The tails are uniformly small in $s$.
On each compact $t$-interval the integrands and their parameter derivatives are
jointly continuous. Apply \KRef{K03.14.improper_derivative} separately at the
two ends, and iterate. The same uniform-tail argument gives continuity of each
derived integral. Since every positive parameter lies in such an interval,
the asserted smoothness follows.
\end{proof}

\KEntry{K03.18.gamma_logconvex}{Log-convexity}
\begin{theorem}
$\log\Gamma$ is convex on $(0,\infty)$.
\end{theorem}
\begin{proof}
For $0<\theta<1$ write
\[
 t^{(1-\theta)p+\theta q-1}e^{-t}
 =(t^{p-1}e^{-t})^{1-\theta}(t^{q-1}e^{-t})^\theta.
\]
Holder's inequality on $[\eta,R]$, with exponents $1/(1-\theta)$ and
$1/\theta$, bounds this integral by
$(\int_\eta^R t^{p-1}e^{-t})^{1-\theta}
 (\int_\eta^R t^{q-1}e^{-t})^\theta$.
All three positive integrals have independent finite endpoint limits, so passage
to those limits gives
$\Gamma((1-\theta)p+\theta q)\leq\Gamma(p)^{1-\theta}\Gamma(q)^\theta$.
Taking logarithms is legitimate by positivity. The cases $\theta=0,1$ are
equalities. This is exactly convexity, using only the properly proved
Riemann integral inequality \KRef{K03.09.holder}.
\end{proof}

\KEntry{K03.18.bohr_mollerup}{The Bohr--Mollerup characterization}
\begin{theorem}
If $F:(0,\infty)\to(0,\infty)$ satisfies $F(1)=1$,
$F(x+1)=xF(x)$ and convexity of $\log F$, then $F=\Gamma$.
\end{theorem}
\begin{proof}
Put $G=\log F$. Recurrence gives $F(n)=(n-1)!$.
For $0<x\leq1$ and integer $n\geq2$, the ordered-secant-slope inequalities
for the convex function $G$ yield
\[
 \log(n-1)\leq\frac{G(n+x)-G(n)}x\leq\log n.
\]
The upper inequality uses $[n,n+x]\subset[n,n+1]$; the lower one compares
with $[n-1,n]$. Since
$F(n+x)=F(x)\prod_{j=0}^{n-1}(x+j)$, exponentiating gives
\[
 \frac{(n-1)!(n-1)^x}{\prod_{j=0}^{n-1}(x+j)}
 \leq F(x)\leq
 \frac{(n-1)!n^x}{\prod_{j=0}^{n-1}(x+j)}.
\]
The left member tends to $\Gamma(x)$ by \eqref{eq:gamma-euler}; the quotient
of the upper and lower members is $(n/(n-1))^x\to1$. Squeeze proves
$F(x)=\Gamma(x)$ on $(0,1]$, and recurrence extends it to all positive $x$.
Gamma itself has all three properties, already proved above. Recurrence and
normalization alone are not unique: multiplication of Gamma by a positive
nonconstant period-one function equal to one at $1$ preserves them.
\end{proof}

\KEntry{K03.18.beta_gamma}{The Beta--Gamma identity without a hidden double integral}
\begin{theorem}
For $p,q>0$,
\[
 B(p,q)=\frac{\Gamma(p)\Gamma(q)}{\Gamma(p+q)}.
\]
\end{theorem}
\begin{proof}
The Beta recurrence gives
$B(p+n,q)=(p)_n B(p,q)/(p+q)_n$, where
$(p)_n=p(p+1)\cdots(p+n-1)$.
We first prove $n^qB(p+n,q)\to\Gamma(q)$.
With $u=n(1-t)$ the left member is
\[
 \int_0^n u^{q-1}(1-u/n)^{n+p-1}\dd u.
\]
For $n\geq2$, $n+p-1\geq n/2$, so this integrand is bounded by
$u^{q-1}e^{-u/2}$. On each $[\eta,R]$ it converges uniformly to
$u^{q-1}e^{-u}$. At zero use the bound $u^{q-1}$, and at infinity use
the displayed exponential bound. The three-interval argument in
\KRef{K03.18.gamma_euler_limit} therefore proves the claimed limit.

The quotient of the Euler product approximants for $\Gamma(p+q)$ and
$\Gamma(p)$, with index $n-1$, shows
\[
 (n-1)^q\frac{(p)_n}{(p+q)_n}\longrightarrow
 \frac{\Gamma(p+q)}{\Gamma(p)}.
\]
Every quantity is positive and the limit is nonzero, so inversion and
$(n/(n-1))^q\to1$ give
$\frac{(p+q)_n}{n^q(p)_n}\to\Gamma(p)/\Gamma(p+q)$.
Multiply this by the established limit of $n^qB(p+n,q)$ and use the exact
finite recurrence identity. This yields the formula without a polar,
Jacobian or full Fubini theorem.
\end{proof}

\KEntry{K03.18.wallis}{Wallis integrals and the product}
\begin{theorem}
Let $I_n=\int_0^{\pi/2}\sin^n x\dd x$, $n\geq0$. Then $I_0=\pi/2$,
$I_1=1$, $I_n=(n-1)I_{n-2}/n$ for $n\geq2$, and
\[
 \prod_{k=1}^n\frac{(2k)^2}{(2k-1)(2k+1)}\longrightarrow\frac\pi2.
\]
Writing $c_n=\binom{2n}{n}/4^n$, we also have
$\sqrt n\,c_n\to1/\sqrt\pi$.
\end{theorem}
\begin{proof}
The first two integrals follow from elementary primitives. Integrating
$\sin^{n-1}x\sin x$ by parts for $n\geq2$ gives
$I_n=(n-1)\int\sin^{n-2}x\cos^2x\dd x=(n-1)(I_{n-2}-I_n)$.
Induction gives
\[
 I_{2n}=\frac\pi2c_n,\qquad I_{2n+1}=\frac1{(2n+1)c_n}.
\]
Since $0\leq\sin x\leq1$,
$I_{2n+2}\leq I_{2n+1}\leq I_{2n}$.
The quotient $I_{2n+2}/I_{2n}=(2n+1)/(2n+2)\to1$, hence
$I_{2n+1}/I_{2n}\to1$. This quotient equals $2W_n/\pi$, where $W_n$
is the displayed finite product, proving Wallis' limit. Further,
$I_{2n}I_{2n+1}=\pi/[2(2n+1)]$. Combined with their quotient limit this
gives $I_{2n}\sim\sqrt\pi/(2\sqrt n)$ and the stated limit for $c_n$.
\end{proof}

\KEntry{K03.18.gaussian}{The Gaussian integral by a one-variable route}
\begin{theorem}
$\Gamma(1/2)=\sqrt\pi$ and $\int_{-\infty}^{\infty}e^{-x^2}\dd x=\sqrt\pi$.
\end{theorem}
\begin{proof}
The Euler approximant at $s=1/2$ simplifies algebraically to
\[
 \frac{n!\sqrt n}{(1/2)(3/2)\cdots(n+1/2)}
 =\frac{2\sqrt n}{(2n+1)c_n}\longrightarrow\sqrt\pi
\]
by Wallis' limit. It therefore equals $\Gamma(1/2)$ in the limit.
Substitution $t=x^2$ on positive regular intervals gives
$\Gamma(1/2)=2\int_0^\infty e^{-x^2}\dd x$, including the endpoint limits.
Evenness gives the two-sided formula. Both tails are absolutely convergent;
there is no principal-value interpretation and no use of a two-dimensional
coordinate substitution.
\end{proof}

\KEntry{K03.18.dirichlet_integral}{The Dirichlet sine integral}
\begin{theorem}
With $\sin x/x$ assigned its continuous value $1$ at zero,
\[
 \int_0^\infty\frac{\sin x}{x}\dd x=\frac\pi2.
\]
\end{theorem}
\begin{proof}
For $a>0$ define $J(a)=\int_0^\infty e^{-ax}\sin x/x\dd x$.
Near zero the integrand is bounded; at infinity it is absolutely integrable.
On a compact parameter interval $a\geq\alpha>0$, the parameter derivative
is $-e^{-ax}\sin x$, bounded by $e^{-\alpha x}$. Thus the proved improper
parameter differentiation rule gives
\[
 J'(a)=-\int_0^\infty e^{-ax}\sin x\dd x=-\frac1{1+a^2}.
\]
The last integral follows by evaluating the primitive
$e^{-ax}(-a\sin x-\cos x)/(1+a^2)$ at zero and infinity.
Also $|J(a)|\leq\int_0^\infty e^{-ax}\dd x=1/a$, since $|\sin x|\leq x$.
The function $\arctan(1/a)$ has the same derivative and tends to zero as
$a\to\infty$; constancy of their difference gives $J(a)=\arctan(1/a)$.

It remains to remove damping, rather than assume it legitimate. For $a\geq0$
and $R>0$, set $g_a(x)=e^{-ax}/x$. It is nonnegative decreasing, with
$g_a(R)\leq1/R$. The primitive $F_R(y)=\int_R^y\sin x\dd x$ has absolute
value at most $2$. Integration by parts gives, for $S>R$,
\[
 \left|\int_R^S g_a(x)\sin x\dd x\right|
 \leq2g_a(S)+2(g_a(R)-g_a(S))\leq\frac2R.
\]
The same bound holds on the infinite tail, uniformly even at $a=0$;
Dirichlet convergence at $a=0$ is included. On $[0,R]$ the damped integrand
converges uniformly to $\sin x/x$ as $a\downarrow0$. Split at $R$, make both
tails small, and then pass to the limit on the compact interval. Therefore
the undamped integral is $\lim_{a\downarrow0}J(a)=\pi/2$.
\end{proof}

\KEntry{K03.18.other_special_functions}{Boundary of the special-function supplement}
\begin{remark}[Routing record, not a new asserted theorem]
Incomplete Gamma, error and Fresnel integrals can be studied after their specific
convergence and error estimates are supplied. General Bessel, Airy, Legendre and
hypergeometric theories belong to D13/D24. Their names here do not claim that
their defining equations, existence theory or asymptotic formulas have been proved
in these notes. No proof above depends on such a theory.
\end{remark}
